% GENERATED FILE. Edit canonical lessons, then run npm run generate:books.
% canonical-source-hash: fnv1a64:9b6f716e04099c9c
% canonical-lessons: TA-C536-kalai-art, TA-C536-oviyam, TA-C536-turikai, TA-C536-tittu-paint, TA-C536-kankatci

\chapter{Painting}
\label{ch:ta-painting}
\hlchaptermodality{536}

\begin{chapteropening}
\textbf{By the end of this chapter:} \emph{I can talk about painting and exhibitions.}

Complete the last lesson of chapter 536: I can talk about painting and exhibitions.
\end{chapteropening}

\section[\texorpdfstring{\ta{கலை} (kalai)}{கலை (kalai)}]{\ta{கலை} (kalai) — an art; the arts}

\label{lesson:TA-C536-kalai-art}



\begin{quote}
Before the new one: say the Tamil for to take aim, then the Tamil for a bat.
\end{quote}

\subsection*{You'll want to know: \ta{கலை}}
\textbf{\ta{கலை}} — \emph{kalai} — \textquotedblleft{}an art; the arts\textquotedblright{}.

\begin{quote}
\ta{நடனம்} \ta{ஒரு} \ta{பழமையான} \ta{கலை}.
\end{quote}

\begin{quote}
\emph{naṭaṉam oru paḻamaiyāṉa kalai.}
\end{quote}

\begin{quote}
Dance is an ancient art.
\end{quote}

An artist is a \ta{கலைஞர்} (\emph{kalaiñar}): \ta{கலை} + -\ta{ஞர்}, \textquotedblleft{}one who knows\textquotedblright{}. The plural is \ta{கலைகள்}.

\begin{cousinweb}
\textbf{\ta{கலை}} (\emph{kalai}) is native Dravidian: DEDR 1297 gives \ta{கல்} (\emph{kal}), \textquotedblleft{}to learn, to practise an art\textquotedblright{}, and beside it \ta{கலை} (\emph{kalai}), \textquotedblleft{}the arts and sciences, learning\textquotedblright{}.
\end{cousinweb}

\begin{grammarlens}[title={What you've built}]
One more word for describing what you lived and felt.
\end{grammarlens}

\subsection*{Your turn}
\begin{itemize}
\raggedright
  \item \emph{Say it:} \emph{kalai}
  \item \emph{Say it:} \emph{kalai}, once more
  \item \emph{Say it:} \emph{naṭaṉam oru paḻamaiyāṉa kalai.} — Dance is an ancient art.
  \item \emph{Recall:} say the Tamil for to take aim, then the Tamil for a bat, then say \emph{kalai} again
\end{itemize}

\begin{tcolorbox}[breakable,colback=teal!4,colframe=teal!35!black,title={Before you move on}]
What does \ta{கலை} mean? (\textquotedblleft{}An art; the arts\textquotedblright{}.) Say it once more.
\end{tcolorbox}
\section[\texorpdfstring{\ta{ஓவியம்} (ōviyam)}{ஓவியம் (ōviyam)}]{\ta{ஓவியம்} (ōviyam) — a painting}

\label{lesson:TA-C536-oviyam}



\begin{quote}
Before the new one: say the Tamil for a bat, then the Tamil for an art.
\end{quote}

\subsection*{You'll want to know: \ta{ஓவியம்}}
\textbf{\ta{ஓவியம்}} — \emph{ōviyam} — \textquotedblleft{}a painting\textquotedblright{}.

\begin{quote}
\ta{அவள்} \ta{வரைந்த} \ta{ஓவியம்} \ta{சுவரில்} \ta{தொங்குகிறது}.
\end{quote}

\begin{quote}
\emph{avaḷ varainta ōviyam cuvaril toṅkukiṟatu.}
\end{quote}

\begin{quote}
The painting she made hangs on the wall.
\end{quote}

A painter is an \ta{ஓவியர்}. A -\ta{அம்} noun: \ta{ஓவியத்தை}, \ta{ஓவியத்தில்}.

\begin{cousinweb}
The origin of \textbf{\ta{ஓவியம்}} (\emph{ōviyam}) is not settled: it is an old literary word, but no dictionary we can check gives its source.
\end{cousinweb}

\begin{grammarlens}[title={What you've built}]
One more word for describing what you lived and felt.
\end{grammarlens}

\subsection*{Your turn}
\begin{itemize}
\raggedright
  \item \emph{Say it:} \emph{ōviyam}
  \item \emph{Say it:} \emph{ōviyam}, once more
  \item \emph{Say it:} \emph{avaḷ varainta ōviyam cuvaril toṅkukiṟatu.} — The painting she made hangs on the wall.
  \item \emph{Recall:} say the Tamil for a bat, then the Tamil for an art, then say \emph{ōviyam} again
\end{itemize}

\begin{tcolorbox}[breakable,colback=teal!4,colframe=teal!35!black,title={Before you move on}]
What does \ta{ஓவியம்} mean? (\textquotedblleft{}A painting\textquotedblright{}.) Say it once more.
\end{tcolorbox}
\section[\texorpdfstring{\ta{தூரிகை} (tūrikai)}{தூரிகை (tūrikai)}]{\ta{தூரிகை} (tūrikai) — a paintbrush}

\label{lesson:TA-C536-turikai}



\begin{quote}
Before the new one: say the Tamil for an art, then the Tamil for a painting.
\end{quote}

\subsection*{You'll want to know: \ta{தூரிகை}}
\textbf{\ta{தூரிகை}} — \emph{tūrikai} — \textquotedblleft{}a paintbrush\textquotedblright{}.

\begin{quote}
\ta{ஓவியர்} \ta{தூரிகையை} \ta{நீரில்} \ta{கழுவினார்}.
\end{quote}

\begin{quote}
\emph{ōviyar tūrikaiyai nīril kaḻuviṉār.}
\end{quote}

\begin{quote}
The painter washed the brush in water.
\end{quote}

A brush for paint; a brush for sweeping is a \ta{துடைப்பம்} (\emph{tuṭaippam}).

\begin{cousinweb}
\textbf{\ta{தூரிகை}} (\emph{tūrikai}) is borrowed from Sanskrit \emph{tūlikā}, \textquotedblleft{}a painter's brush, a pencil\textquotedblright{} (Monier-Williams).
\end{cousinweb}

\begin{grammarlens}[title={What you've built}]
One more word for describing what you lived and felt.
\end{grammarlens}

\subsection*{Your turn}
\begin{itemize}
\raggedright
  \item \emph{Say it:} \emph{tūrikai}
  \item \emph{Say it:} \emph{tūrikai}, once more
  \item \emph{Say it:} \emph{ōviyar tūrikaiyai nīril kaḻuviṉār.} — The painter washed the brush in water.
  \item \emph{Recall:} say the Tamil for an art, then the Tamil for a painting, then say \emph{tūrikai} again
\end{itemize}

\begin{tcolorbox}[breakable,colback=teal!4,colframe=teal!35!black,title={Before you move on}]
What does \ta{தூரிகை} mean? (\textquotedblleft{}A paintbrush\textquotedblright{}.) Say it once more.
\end{tcolorbox}
\section[\texorpdfstring{\ta{தீட்டு} (tīṭṭu)}{தீட்டு (tīṭṭu)}]{\ta{தீட்டு} (tīṭṭu) — to paint a picture; to sharpen}

\label{lesson:TA-C536-tittu-paint}



\begin{quote}
Before the new one: say the Tamil for a painting, then the Tamil for a paintbrush.
\end{quote}

\subsection*{You'll want to know: \ta{தீட்டு}}
\textbf{\ta{தீட்டு}} — \emph{tīṭṭu} — \textquotedblleft{}to paint a picture; to sharpen\textquotedblright{}.

\begin{quote}
\ta{அவன்} \ta{சுவரில்} \ta{ஒரு} \ta{பெரிய} \ta{மயிலைத்} \ta{தீட்டினான்}.
\end{quote}

\begin{quote}
\emph{avaṉ cuvaril oru periya mayilait tīṭṭiṉāṉ.}
\end{quote}

\begin{quote}
He painted a big peacock on the wall.
\end{quote}

Past (I …): \ta{தீட்டினேன்}, \emph{tīṭṭiṉēṉ}.

Two senses from one verb: \ta{படம்} \ta{தீட்டு}, to paint a picture, and \ta{கத்தியைத்} \ta{தீட்டு}, to sharpen a knife.

\begin{cousinweb}
\textbf{\ta{தீட்டு}} (\emph{tīṭṭu}) is native Dravidian: DEDR 3273 gives \ta{தீட்டு} (\emph{tīṭṭu}), \textquotedblleft{}to whet, to sharpen, to polish, to smear, to draw\textquotedblright{}.
\end{cousinweb}

\begin{grammarlens}[title={What you've built}]
One more word for describing what you lived and felt.
\end{grammarlens}

\subsection*{Your turn}
\begin{itemize}
\raggedright
  \item \emph{Say it:} \emph{tīṭṭu}
  \item \emph{Say it:} \emph{tīṭṭu}, once more
  \item \emph{Say it:} \emph{avaṉ cuvaril oru periya mayilait tīṭṭiṉāṉ.} — He painted a big peacock on the wall.
  \item \emph{Recall:} say the Tamil for a painting, then the Tamil for a paintbrush, then say \emph{tīṭṭu} again
\end{itemize}

\begin{tcolorbox}[breakable,colback=teal!4,colframe=teal!35!black,title={Before you move on}]
What does \ta{தீட்டு} mean? (\textquotedblleft{}To paint a picture; to sharpen\textquotedblright{}.) Say it once more.
\end{tcolorbox}
\section[\texorpdfstring{\ta{கண்காட்சி} (kaṇkāṭci)}{கண்காட்சி (kaṇkāṭci)}]{\ta{கண்காட்சி} (kaṇkāṭci) — an exhibition}

\label{lesson:TA-C536-kankatci}



\begin{quote}
Before the new one: say the Tamil for a paintbrush, then the Tamil for to paint a picture.
\end{quote}

\subsection*{You'll want to know: \ta{கண்காட்சி}}
\textbf{\ta{கண்காட்சி}} — \emph{kaṇkāṭci} — \textquotedblleft{}an exhibition\textquotedblright{}.

\begin{quote}
\ta{நகர} \ta{மண்டபத்தில்} \ta{ஒரு} \ta{புகைப்படக்} \ta{கண்காட்சி} \ta{நடக்கிறது}.
\end{quote}

\begin{quote}
\emph{nakara maṇṭapattil oru pukaippaṭak kaṇkāṭci naṭakkiṟatu.}
\end{quote}

\begin{quote}
A photography exhibition is on in the town hall.
\end{quote}

Word for word \textquotedblleft{}an eye-sight\textquotedblright{}: \ta{கண்} + \ta{காட்சி} \textquotedblleft{}a sight\textquotedblright{}, from \ta{காண்} (\emph{kāṇ}), to see.

\begin{cousinweb}
\textbf{\ta{கண்காட்சி}} (\emph{kaṇkāṭci}) is a native compound: \ta{கண்} (\emph{kaṇ}), \textquotedblleft{}an eye\textquotedblright{} (DEDR 1159), and \ta{காட்சி} (\emph{kāṭci}), \textquotedblleft{}a sight, a show\textquotedblright{}, from \ta{காண்} (\emph{kāṇ}), \textquotedblleft{}to see\textquotedblright{} (DEDR 1443).
\end{cousinweb}

\begin{grammarlens}[title={What you've built}]
One more word for describing what you lived and felt.
\end{grammarlens}

\subsection*{Your turn}
\begin{itemize}
\raggedright
  \item \emph{Say it:} \emph{kaṇkāṭci}
  \item \emph{Say it:} \emph{kaṇkāṭci}, once more
  \item \emph{Say it:} \emph{nakara maṇṭapattil oru pukaippaṭak kaṇkāṭci naṭakkiṟatu.} — A photography exhibition is on in the town hall.
  \item \emph{Recall:} say the Tamil for a paintbrush, then the Tamil for to paint a picture, then say \emph{kaṇkāṭci} again
\end{itemize}

\begin{tcolorbox}[breakable,colback=teal!4,colframe=teal!35!black,title={Before you move on}]
What does \ta{கண்காட்சி} mean? (\textquotedblleft{}An exhibition\textquotedblright{}.) Say it once more.
\end{tcolorbox}
